% ── 附录 ──
% \label{sec:appendix} 不得删除：latex_gate.py 依赖它切分正文页数与附录页数。
% 附录承载支撑材料索引、核心代码、环境和复现说明。

\makeatletter
\let\@oldaddcontentsline\addcontentsline
\renewcommand{\addcontentsline}[3]{}
\section*{附录}
\label{sec:appendix}
\let\addcontentsline\@oldaddcontentsline
\makeatother

\subsection*{附录1：支撑材料目录}
\label{app:support-files}

逐项列出随论文提交的支撑材料，路径与 delivery-manifest.json 保持一致。竞赛已提供的原始数据不重复提交；自行获得的数据才可在支撑材料中列出。

\begin{table}[H]
  \centering
  \caption{支撑材料目录}
  \label{tab:appendix_files}
  \renewcommand{\arraystretch}{1.2}
  \begin{tabularx}{\textwidth}{LXX}
    \toprule
    相对路径 & 类型 & 用途与论文位置 \\
    \midrule
    支撑材料/source/ & 完整源码 & 完整、可运行的项目；入口为支撑材料/source/main.py。 \\
    支撑材料/data/ & 自行获得数据 & 仅按实际需要提交，并写明来源与处理方式。 \\
    支撑材料/results/ & 补充结果 & 放正文容纳不下的结果、表格或图表。 \\
    \bottomrule
  \end{tabularx}
\end{table}

\subsection*{附录2：核心源程序节选}
\label{app:core-code}

本节只放说明模型逻辑所需的核心代码，不打印整个项目。完整可运行源码及依赖文件位于支撑材料/source/，实际入口为支撑材料/source/main.py。

\begin{verbatim}
# Replace with the model's core logic.
def solve(instance):
    # model construction, key constraint, and decision logic
    ...
\end{verbatim}

节选应对应正文的模型、约束或求解逻辑；不得把这一小段节选说成完整源码。

\subsection*{附录3：运行环境与复现说明}
\label{app:reproduction}

请如实写明语言、版本、依赖、操作系统条件、输入输出位置和复现命令。最终包根目录的 reproduce.py 应执行至少一条真实的最小重算命令，而不是只做文件存在性检查。

\begin{verbatim}
# Example only. Replace with the actual environment and command.
python --version
python reproduce.py
\end{verbatim}
